\section{Transformer Encoder Architecture}\label{sec:method:arch}

The encoder $\Enc$ takes the form of a transformer-encoder backbone \parencite{vaswani2017attention} with a shared trunk followed by two task-specific branches for mode clustering and deinterleaving. The overall layout is shown in \cref{fig:method:arch}.
The trunk produces a single contextualised representation of the input window, which the branches specialise into emitter-oriented and mode-oriented embeddings for the two estimation tasks defined in \cref{sec:method:formulation}.

\begin{figure}[tbp]
  \centering
  % =====================================================================
% Schematic of the Y-shaped transformer encoder f_theta used in
% \cref{sec:method:arch}.  Designed for the single-column thesis body
% (text width ~155 mm); portrait orientation, bottom-up data flow.
% All affine projections (input embedding, branch projections, and
% projection heads) share the same visual style, since they are the
% same kind of layer.
% Loaded from sections/04_method/03_architecture.tex via \input{...}.
% =====================================================================
\begin{tikzpicture}[
    font=\small,
    >=Latex,
    node distance=6mm and 8mm,
    every node/.style={transform shape},
    block/.style={
      draw, thick, rounded corners=2pt,
      align=center,
      inner sep=4pt, inner ysep=5pt,
      minimum height=9mm,
    },
    io/.style={align=center, font=\small\itshape, inner sep=2pt},
    affine/.style={block, fill=orange!12, minimum width=42mm},
    trunk/.style={block,  fill=gray!15,   minimum width=62mm},
    branch/.style={block, fill=blue!15,   minimum width=42mm},
    norm/.style={block,   fill=green!15,  minimum width=42mm,
                 minimum height=6mm},
    arrow/.style={->, thick, shorten <=1pt, shorten >=1pt},
    line/.style={thick,  shorten <=1pt, shorten >=1pt},
  ]

  \node[io] (in)
    {\glspl{pdw}: $\vect{x}_1,\ldots,\vect{x}_L \in \R^{d}$\quad ($d{=}7$)};

  \node[affine, minimum width=62mm, above=of in] (embed)
    {Input embedding\\
     $\vect{h}^{(0)}_n = \mat{W}_{\mathrm{in}}\vect{x}_n + \vect{b}_{\mathrm{in}}$
     \;:\; $\R^{d}\!\to\!\R^{256}$};

  \node[trunk, above=of embed] (trunk)
    {Shared transformer trunk\\
     $N_{\mathrm{sh}}{=}4$ encoder layers\\
     $d_{\mathrm{model}}{=}256$,\; $h{=}8$,\; $d_{\mathrm{ff}}{=}2048$};

  \coordinate (fork) at ($(trunk.north)+(0,5mm)$);

  \node[affine] (proj_em) at ($(fork)+(-32mm, 9mm)$)
    {Affine projection\\$\R^{256}\!\to\!\R^{256}$};
  \node[affine] (proj_md) at ($(fork)+( 32mm, 9mm)$)
    {Affine projection\\$\R^{256}\!\to\!\R^{256}$};

  \node[branch, above=of proj_em] (br_em)
    {Emitter branch\\
     $N_{\mathrm{br}}{=}4$ encoder layers\\
     $d_{\mathrm{br}}{=}256$,\; $h{=}8$,\; $d_{\mathrm{ff}}{=}2048$};
  \node[branch, above=of proj_md] (br_md)
    {Mode branch\\
     $N_{\mathrm{br}}{=}4$ encoder layers\\
     $d_{\mathrm{br}}{=}256$,\; $h{=}8$,\; $d_{\mathrm{ff}}{=}2048$};

  \node[affine, above=of br_em] (head_em)
    {Projection head\\
     $\mat{W}_{\mathrm{em}}\vect{h}^{\mathrm{em}}_n + \vect{b}_{\mathrm{em}}$\\
     $\R^{256}\!\to\!\R^{64}$};
  \node[affine, above=of br_md] (head_md)
    {Projection head\\
     $\mat{W}_{\mathrm{md}}\vect{h}^{\mathrm{md}}_n + \vect{b}_{\mathrm{md}}$\\
     $\R^{256}\!\to\!\R^{64}$};

  \node[norm, above=of head_em] (norm_em)
    {$\ell_2$ normalization};
  \node[norm, above=of head_md] (norm_md)
    {$\ell_2$ normalization};

  \node[io, above=of norm_em] (out_em)
    {Emitter embedding\\$\vect{z}^{\mathrm{em}}_n \in \R^{64}$};
  \node[io, above=of norm_md] (out_md)
    {Mode embedding\\$\vect{z}^{\mathrm{md}}_n \in \R^{64}$};

  \draw[arrow] (in)        -- (embed);
  \draw[arrow] (embed)     -- (trunk);
  \draw[line]  (trunk.north) -- (fork);
  \draw[arrow] (fork) -| (proj_em.south);
  \draw[arrow] (fork) -| (proj_md.south);
  \draw[arrow] (proj_em)   -- (br_em);
  \draw[arrow] (proj_md)   -- (br_md);
  \draw[arrow] (br_em)     -- (head_em);
  \draw[arrow] (br_md)     -- (head_md);
  \draw[arrow] (head_em)   -- (norm_em);
  \draw[arrow] (head_md)   -- (norm_md);
  \draw[arrow] (norm_em)   -- (out_em);
  \draw[arrow] (norm_md)   -- (out_md);

\end{tikzpicture}

  \caption{Architecture of the encoder $\Enc$. A sequence of \glspl{pdw} is first lifted to $d_{\mathrm{model}}=256$ by an affine input embedding, processed by a shared trunk of $6$ transformer-encoder layers, and then split into two parallel branches, each with its own affine $256\!\to\!256$ projection and $2$ encoder layers. Branch outputs are mapped by affine projection heads to the $64$-dimensional emitter and mode embedding spaces and $\ell_2$-normalized over the embedding dimension. The input embedding, branch projections, and projection heads are all affine maps of the same form.}
  \label{fig:method:arch}
\end{figure}

\subsection*{Input embedding}

Each pulse vector $\vect{x}_n \in \R^{d}$, built from the scaled \glspl{pdw} of \cref{sec:method:data}, is mapped to a token embedding through an affine projection
\begin{equation}\label{eq:method:input_embed}
  \vect{h}^{(0)}_n = \mat{W}_{\mathrm{in}}\vect{x}_n + \vect{b}_{\mathrm{in}},
  \qquad
  \vect{h}^{(0)}_n \in \R^{d_{\mathrm{model}}},
\end{equation}
where $\mat{W}_{\mathrm{in}} \in \R^{d_{\mathrm{model}} \times d}$ and $d_{\mathrm{model}}$ is the trunk width specified below.
In the baseline configuration of the encoder, no positional encoding is added to $\vect{h}^{(0)}_n$.
The temporal ordering of pulses is already carried by the \gls{toa} entry of $\vect{x}_n$; moreover, the index-based absolute positional encodings of the original architecture \parencite{vaswani2017attention} are ill-suited to this input, since the \gls{toa} gap between consecutive pulses is not constant and an ordinal position therefore misrepresents the relative distances in time.
Whether an explicit \emph{relative} encoding of time and arrival direction nevertheless adds information beyond their presence as input features is revisited in \cref{sec:method:attention}, where a pairwise attention bias and a rotary positional encoding acting on physical coordinate differences are evaluated as drop-in modifications of this architecture.

\subsection*{Shared trunk}

The token sequence $(\vect{h}^{(0)}_1,\ldots,\vect{h}^{(0)}_L)$ is processed by a stack of $N_{\mathrm{sh}}$ standard transformer-encoder layers, each comprising a \gls{mha} sub-layer and a position-wise \gls{ffn} sub-layer with residual connections and layer normalisation around both.
Self-attention is unmasked, so every pulse can attend to every other pulse in the window.
The trunk operates with hidden dimension $d_{\mathrm{model}} = 256$, $h_{\mathrm{sh}} = 8$ attention heads (so that each head has dimension $32$), and an inner \gls{ffn} dimension of $d_{\mathrm{ff,sh}} = 2048$; $N_{\mathrm{sh}} = 6$ layers are used.
Its output is a contextualised token sequence $\mat{H}^{\mathrm{sh}} = (\vect{h}^{\mathrm{sh}}_1,\ldots,\vect{h}^{\mathrm{sh}}_L)$ with $\vect{h}^{\mathrm{sh}}_n \in \R^{256}$, in which information has been mixed across the full window.

\subsection*{Task-specific branches}

The trunk output feeds two parallel branches that do not share parameters.
Each branch begins with an affine projection of the same form as \cref{eq:method:input_embed}, after which $N_{\mathrm{br}} = 2$ transformer-encoder layers of the same form as the trunk are applied.
The projection decouples the branch representation from the shared trunk output and, more generally, allows the branch width to differ from the trunk width; here the branches retain $d_{\mathrm{br}} = 256$, so the projection maps $\R^{256} \to \R^{256}$.
Within a branch, \gls{mha} uses $h_{\mathrm{br}} = 4$ heads (per-head dimension $64$) and the \gls{ffn} inner dimension is $d_{\mathrm{ff,br}} = 1024$.
The reduced depth and \gls{ffn} width reflect a deliberate asymmetry: the trunk is intended to carry the heavy lifting of contextual mixing, while the branches only need to specialise the representation to either emitter identity or operating mode.
The larger per-head dimension in the branches also anticipates the multi-coordinate rotary encoding of \cref{sec:method:attention}, which partitions the rotation planes of a head across three coordinates and therefore benefits from more planes per head; keeping the same dimensions in the baseline ensures that the attention variants are compared on an otherwise identical architecture.
Write the resulting token sequences as $\mat{H}^{\mathrm{em}} = (\vect{h}^{\mathrm{em}}_1,\ldots,\vect{h}^{\mathrm{em}}_L)$ and $\mat{H}^{\mathrm{md}} = (\vect{h}^{\mathrm{md}}_1,\ldots,\vect{h}^{\mathrm{md}}_L)$ with $\vect{h}^{\mathrm{em}}_n, \vect{h}^{\mathrm{md}}_n \in \R^{256}$.

\subsection*{Projection heads}

Each branch terminates in an affine projection head---of the same form as the input embedding \cref{eq:method:input_embed}---that maps every token to its respective embedding space, followed by $\ell_2$ normalization over the embedding dimension,
\begin{equation}\label{eq:method:proj_heads}
  \vect{z}^{\mathrm{em}}_n
  = \frac{\mat{W}_{\mathrm{em}}\vect{h}^{\mathrm{em}}_n + \vect{b}_{\mathrm{em}}}
         {\bigl\lVert \mat{W}_{\mathrm{em}}\vect{h}^{\mathrm{em}}_n + \vect{b}_{\mathrm{em}} \bigr\rVert_2},
  \qquad
  \vect{z}^{\mathrm{md}}_n
  = \frac{\mat{W}_{\mathrm{md}}\vect{h}^{\mathrm{md}}_n + \vect{b}_{\mathrm{md}}}
         {\bigl\lVert \mat{W}_{\mathrm{md}}\vect{h}^{\mathrm{md}}_n + \vect{b}_{\mathrm{md}} \bigr\rVert_2},
\end{equation}
with $\vect{z}^{\mathrm{em}}_n, \vect{z}^{\mathrm{md}}_n \in \R^{p}$ for a common embedding dimension $p = 64$.
Collecting the per-pulse outputs gives the two embedding sequences
\begin{equation}\label{eq:method:dual_embeddings}
  \mat{Z}^{\mathrm{em}} = (\vect{z}^{\mathrm{em}}_1,\ldots,\vect{z}^{\mathrm{em}}_L),
  \qquad
  \mat{Z}^{\mathrm{md}} = (\vect{z}^{\mathrm{md}}_1,\ldots,\vect{z}^{\mathrm{md}}_L).
\end{equation}
The normalization constrains both embedding sequences to the unit sphere in $\R^{64}$, so cluster structure is expressed through angular separation rather than vector magnitude; this is the geometry assumed by the cosine-similarity contrastive objective of \cref{sec:method:loss}.

\subsection*{From continuous embeddings to discrete partitions}

The projection heads in \cref{eq:method:proj_heads} produce continuous vectors, whereas the per-pulse outputs in \cref{eq:method:estimated_label_pair} are discrete.
At inference time, each branch's $L$-token embedding sequence is therefore decoded independently and per window by a fixed clustering operator,
\begin{equation}\label{eq:method:cluster_decoding}
  (\hat{e}_1,\ldots,\hat{e}_L) = \mathrm{DBSCAN}\!\bigl(\{\vect{z}^{\mathrm{em}}_n\}_{n=1}^{L}\bigr),
  \quad
  (\hat{m}_1,\ldots,\hat{m}_L) = \mathrm{DBSCAN}\!\bigl(\{\vect{z}^{\mathrm{md}}_n\}_{n=1}^{L}\bigr),
\end{equation}
with cardinalities $\hat{K},\hat{M} \in \N$ inferred from the data (\cref{fig:method:cluster_decoding}).
\Gls{dbscan} \parencite{ester1996density} groups embeddings by local density at a Euclidean radius $\varepsilon$ and minimum neighbourhood count $\mathrm{minPts}$, assigning sparser points to a noise category.
The radius $\varepsilon$ is selected by a grid search over $\varepsilon \in \{0.1, 0.2, 0.3, 0.4, 0.5\}$ on the validation set, separately for each branch, using the clustering metrics of \cref{sec:method:metrics}; the selected values are reported in \cref{ch:experiments}.
The neighbourhood count is fixed at $\mathrm{minPts} = 5$, following its use in prior work on embedding-based deinterleaving \parencite{gunn2025deinterleaving}.
With these settings the noise category remained marginal on the validation set---empty for most windows and almost never exceeding ten pulses per window---so no further treatment of noise-labelled pulses was deemed necessary.
Clustering per window matches the recording-local supervision scope of the emitter branch and keeps the runtime linear in window length; cross-window alignment of mode embeddings is instead supplied by the batch-wide contrastive term of \cref{sec:method:loss} and resolved by the permutation-invariant metrics of \cref{sec:method:metrics}.
The operator contributes no gradients, and the trunk, branch, and projection-head weights together form the parameter vector $\bm{\theta}$ optimised by the loss in \cref{sec:method:loss}.

\begin{figure}[tbp]
  \centering
  % =====================================================================
% Schematic of the inference-time mapping from continuous branch
% embeddings to discrete per-pulse labels.  Two parallel pipelines
% (emitter / mode) take the $L$-token embedding sequences produced by
% the projection heads of \cref{sec:method:arch} and cluster them
% independently with \gls{hdbscan} on a per-window basis, yielding the
% discrete partitions of
% \cref{eq:method:emitter_predictor,eq:method:mode_predictor}.  The
% pipelines are structurally identical; the asymmetry between the two
% branches lives in the training loss (within-window vs. batch-wide
% supervised contrastive) and in the semantic load of the integer
% indices, not in the inference pipeline.  Loaded from
% sections/04_method/03_architecture.tex via \input{...}.
% =====================================================================
\begin{tikzpicture}[
    font=\small,
    >=Latex,
    block/.style={
      draw, thick, rounded corners=2pt,
      align=center, inner sep=4pt,
      minimum height=8mm, minimum width=48mm,
    },
    scatterbox/.style={
      draw, thick, rounded corners=2pt,
      fill=gray!4,
      minimum width=48mm, minimum height=32mm,
    },
    hdb/.style={block, fill=blue!12},
    outblock/.style={block, fill=green!15},
    io/.style={align=center, font=\small\itshape, inner sep=2pt},
    pt/.style={circle, fill, inner sep=0pt, minimum size=1.8pt},
    arrow/.style={->, thick, shorten <=1pt, shorten >=1pt},
    contour/.style={dashed, thick},
  ]

  % =================== EMITTER COLUMN ===================
  \begin{scope}[xshift=0mm]
    \node[font=\footnotesize\bfseries, anchor=south]
      at (0, 3mm) {Emitter branch};

    \node[io] (em_in) at (0, -3mm)
      {Emitter embeddings $\{\vect{z}^{\mathrm{em}}_n\}_{n=1}^{L} \subset \R^{p}$};

    \node[scatterbox] (em_scatter) at (0, -22mm) {};

    % Cluster A
    \draw[contour, color=blue!70!black]
      ($(em_scatter.center)+(-13mm, 4mm)$) ellipse [x radius=6mm, y radius=4mm];
    \foreach \x/\y in {-15/5, -13/6, -12/3, -10/5, -14/2, -11/4}
      \node[pt, color=blue!70!black] at ($(em_scatter.center)+(\x mm, \y mm)$) {};
    \node[font=\scriptsize, color=blue!70!black]
      at ($(em_scatter.center)+(-13mm, 10mm)$) {$\hat{e}{=}1$};

    % Cluster B
    \draw[contour, color=red!70!black]
      ($(em_scatter.center)+(11mm, 3mm)$) ellipse [x radius=6mm, y radius=4mm];
    \foreach \x/\y in {8/2, 10/4, 12/2, 14/3, 11/1, 13/4}
      \node[pt, color=red!70!black] at ($(em_scatter.center)+(\x mm, \y mm)$) {};
    \node[font=\scriptsize, color=red!70!black]
      at ($(em_scatter.center)+(11mm, 9mm)$) {$\hat{e}{=}2$};

    % Cluster C
    \draw[contour, color=green!45!black]
      ($(em_scatter.center)+(-2mm, -8mm)$) ellipse [x radius=7mm, y radius=4mm];
    \foreach \x/\y in {-5/-7, -3/-9, 0/-8, 2/-7, -1/-10, 3/-9}
      \node[pt, color=green!45!black] at ($(em_scatter.center)+(\x mm, \y mm)$) {};
    \node[font=\scriptsize, color=green!45!black]
      at ($(em_scatter.center)+(-2mm, -13mm)$) {$\hat{e}{=}3$};

    \node[hdb] (em_hdb) at (0, -45mm)
      {\gls{hdbscan} per window};
    \node[outblock] (em_out) at (0, -60mm)
      {Train-local emitter IDs\\$\hat{e}_n \in \{1,\ldots,\hat{K}_{\mathrm{em}}\}$};

    \draw[arrow] (em_in) -- (em_scatter);
    \draw[arrow] (em_scatter) -- (em_hdb);
    \draw[arrow] (em_hdb) -- (em_out);
  \end{scope}

  % =================== MODE COLUMN ===================
  \begin{scope}[xshift=65mm]
    \node[font=\footnotesize\bfseries, anchor=south]
      at (0, 3mm) {Mode branch};

    \node[io] (md_in) at (0, -3mm)
      {Mode embeddings $\{\vect{z}^{\mathrm{md}}_n\}_{n=1}^{L} \subset \R^{q}$};

    \node[scatterbox] (md_scatter) at (0, -22mm) {};

    % Cluster A
    \draw[contour, color=blue!70!black]
      ($(md_scatter.center)+(-13mm, 4mm)$) ellipse [x radius=6mm, y radius=4mm];
    \foreach \x/\y in {-15/5, -13/6, -12/3, -10/5, -14/2, -11/4}
      \node[pt, color=blue!70!black] at ($(md_scatter.center)+(\x mm, \y mm)$) {};
    \node[font=\scriptsize, color=blue!70!black]
      at ($(md_scatter.center)+(-13mm, 10mm)$) {$\hat{m}{=}1$};

    % Cluster B
    \draw[contour, color=red!70!black]
      ($(md_scatter.center)+(11mm, 3mm)$) ellipse [x radius=6mm, y radius=4mm];
    \foreach \x/\y in {8/2, 10/4, 12/2, 14/3, 11/1, 13/4}
      \node[pt, color=red!70!black] at ($(md_scatter.center)+(\x mm, \y mm)$) {};
    \node[font=\scriptsize, color=red!70!black]
      at ($(md_scatter.center)+(11mm, 9mm)$) {$\hat{m}{=}2$};

    % Cluster C
    \draw[contour, color=green!45!black]
      ($(md_scatter.center)+(-2mm, -8mm)$) ellipse [x radius=7mm, y radius=4mm];
    \foreach \x/\y in {-5/-7, -3/-9, 0/-8, 2/-7, -1/-10, 3/-9}
      \node[pt, color=green!45!black] at ($(md_scatter.center)+(\x mm, \y mm)$) {};
    \node[font=\scriptsize, color=green!45!black]
      at ($(md_scatter.center)+(-2mm, -13mm)$) {$\hat{m}{=}3$};

    \node[hdb] (md_hdb) at (0, -45mm)
      {\gls{hdbscan} per window};
    \node[outblock] (md_out) at (0, -60mm)
      {Per-window mode cluster IDs\\$\hat{m}_n \in \{1,\ldots,\hat{M}\}$};

    \draw[arrow] (md_in) -- (md_scatter);
    \draw[arrow] (md_scatter) -- (md_hdb);
    \draw[arrow] (md_hdb) -- (md_out);
  \end{scope}

\end{tikzpicture}

  \caption{Inference-time decoding of branch embeddings into per-pulse cluster indices via \cref{eq:method:cluster_decoding}. Dots are pulse embeddings and dashed ellipses are \gls{dbscan} clusters in one window. The pipeline is identical on both branches and the integer outputs are local to each clustering instance; they acquire their interpretation from the supervision scope of \cref{sec:method:formulation}, which differs between the two branches.}
  \label{fig:method:cluster_decoding}
\end{figure}
